% ============================================================
% DOCUMENTCLASS
% ============================================================
\documentclass[12pt,oneside]{book}

% ============================================================
% METADATOS
% ============================================================
\newcommand{\universidad}{Universidad de Buenos Aires (UBA)}
\newcommand{\facultad}{Facultad de Derecho}
\newcommand{\catedra}{Cátedra de Derechos Reales}
\newcommand{\materia}{Derechos Reales}
\newcommand{\titulo}{Derecho Real de Superficie}
\newcommand{\autor}{Isaac Edgar Camacho Ocampo}
\newcommand{\anio}{2026}

% ============================================================
% PAQUETES
% ============================================================
\usepackage[utf8]{inputenc}
\usepackage[T1]{fontenc}
\usepackage[spanish,es-nodecimaldot]{babel}
\usepackage[
    top=2.5cm,
    bottom=2.5cm,
    left=3cm,
    right=2.5cm,
    headheight=15pt
]{geometry}
\usepackage{amsmath}
\usepackage{amssymb}
\usepackage{mathtools}
\usepackage{graphicx}
\usepackage{float}
\usepackage{caption}
\usepackage{subcaption}
\usepackage{booktabs}
\usepackage{array}
\usepackage{longtable}
\usepackage{tabularx}
\usepackage{enumitem}
\usepackage{xcolor}
\usepackage[most]{tcolorbox}
\usepackage{fancyhdr}
\usepackage{ragged2e}
\usepackage[
    colorlinks=true,
    linkcolor=blue!50!black,
    citecolor=green!40!black,
    urlcolor=blue!60!black,
    pdftitle={\titulo},
    pdfauthor={\autor},
    pdfsubject={Apuntes universitarios},
    bookmarks=true
]{hyperref}

% ============================================================
% CONFIGURACIÓN
% ============================================================
\graphicspath{
    {./img/}
    {./graficos/}
    {./figuras/}
}
\setcounter{tocdepth}{2}
\setlength{\parindent}{0pt}
\setlength{\parskip}{0.5em}
\setlist[itemize]{
    topsep=0.3em,
    itemsep=0.2em
}
\setlist[enumerate]{
    topsep=0.3em,
    itemsep=0.2em
}
\clubpenalty=10000
\widowpenalty=10000

% ============================================================
% COMANDOS
% ============================================================
\newcommand{\abs}[1]{\left|#1\right|}
\newcommand{\norm}[1]{\left\lVert#1\right\rVert}
\newcommand{\vect}[1]{\mathbf{#1}}
\newcommand{\deriv}[2]{\frac{d#1}{d#2}}
\newcommand{\pderiv}[2]{\frac{\partial#1}{\partial#2}}

% ============================================================
% ENTORNOS / CAJAS
% ============================================================
\newtcolorbox{definition}[1]{
    enhanced,
    breakable,
    title={Definición: #1},
    fonttitle=\bfseries,
    colback=blue!3,
    colframe=blue!50!black,
    boxrule=0.8pt,
    arc=2mm
}

\newtcolorbox{important}{
    enhanced,
    breakable,
    title={Importante},
    fonttitle=\bfseries,
    colback=yellow!8,
    colframe=orange!70!black,
    boxrule=0.8pt,
    arc=2mm
}

\newtcolorbox{example}[1]{
    enhanced,
    breakable,
    title={Ejemplo: #1},
    fonttitle=\bfseries,
    colback=green!4,
    colframe=green!40!black,
    boxrule=0.8pt,
    arc=2mm
}

\newtcolorbox{formula}[1]{
    enhanced,
    breakable,
    title={Fórmula / Regla: #1},
    fonttitle=\bfseries,
    colback=purple!4,
    colframe=purple!50!black,
    boxrule=0.8pt,
    arc=2mm
}

\newtcolorbox{exercise}[1]{
    enhanced,
    breakable,
    title={Ejercicio: #1},
    fonttitle=\bfseries,
    colback=gray!5,
    colframe=gray!60!black,
    boxrule=0.8pt,
    arc=2mm
}

\newtcolorbox{note}{
    enhanced,
    breakable,
    title={Nota},
    fonttitle=\bfseries,
    colback=white,
    colframe=gray!50,
    boxrule=0.6pt,
    arc=2mm
}

% ============================================================
% CONFIGURACIÓN FINAL
% ============================================================
\pagestyle{fancy}
\fancyhf{}
\fancyhead[L]{\nouppercase{\leftmark}}
\fancyhead[R]{\materia}
\fancyfoot[C]{\thepage}
\renewcommand{\headrulewidth}{0.4pt}
\renewcommand{\footrulewidth}{0pt}

\fancypagestyle{plain}{
    \fancyhf{}
    \fancyfoot[C]{\thepage}
    \renewcommand{\headrulewidth}{0pt}
}

% ============================================================
% DOCUMENT
% ============================================================
\begin{document}

% ------------------------------------------------------------
% PORTADA
% ------------------------------------------------------------
\begin{titlepage}

\thispagestyle{empty}

\begin{center}

\vspace*{1cm}

{\large \textsc{\universidad}}

\vspace{0.2cm}

{\large \textsc{\facultad}}

\vspace{0.2cm}

{\large \catedra}

\vspace{3cm}

{\Huge \textbf{\materia}}

\vspace{1cm}

{\LARGE \titulo}

\vspace{0.8cm}

\textit{Comprender cada concepto con paciencia es el camino más firme hacia un conocimiento duradero.}

\vspace{1.2cm}

\rule{0.70\textwidth}{0.5pt}

\vspace{0.5cm}

{\large Apuntes de estudio}

\vspace{0.5cm}

\rule{0.70\textwidth}{0.5pt}

\vfill

{\large \autor}

\vspace{0.3cm}

{\large \anio}

\end{center}

\vfill

\begin{flushleft}
\textit{Este es un modesto aporte a los estudiantes de ``Derechos Reales'' no pretende sustituir el material oficial de la materia.}
\end{flushleft}

\end{titlepage}

% ------------------------------------------------------------
% ÍNDICE
% ------------------------------------------------------------
\tableofcontents

% ------------------------------------------------------------
% CONTENIDO
% ------------------------------------------------------------

\chapter[Fundamentos generales]{Fundamentos generales del derecho real de superficie}

\section{Relevancia profesional del tema}

El derecho real de superficie es una de las figuras más innovadoras del Código Civil y Comercial de la Nación (CCCN), vigente desde 2015. Su comprensión habilita a:

\begin{itemize}
    \item Redactar escrituras públicas de constitución de superficie sobre inmuebles urbanos, rurales o forestales.
    \item Asesorar a desarrolladores inmobiliarios que deseen construir en terrenos ajenos sin comprarlos (modelo BTS, \textit{Build-to-Suit}).
    \item Estructurar fideicomisos o emprendimientos de propiedad horizontal levantados sobre suelo de terceros, algo cada vez más frecuente en la Ciudad de Buenos Aires.
    \item Registrar y oponer \textit{erga omnes} el derecho ante el Registro de la Propiedad Inmueble.
    \item Calcular y negociar la indemnización al superficiario cuando el derecho se extingue (Art.~2126 CCCN).
\end{itemize}

\begin{important}
El derecho de superficie \textbf{separa la propiedad del suelo de la propiedad de lo construido, plantado o forestado}. Esta disociación genera oportunidades de negocio que antes eran jurídicamente imposibles en Argentina, y es una herramienta cotidiana en mercados inmobiliarios sofisticados de Europa y los Estados Unidos.
\end{important}

\begin{note}
Como analogía, el derecho de superficie se asemeja a la situación de un inquilino que coloca muebles propios en una casa ajena, con la diferencia de que, en este caso, el «mueble» puede ser un edificio entero.
\end{note}

\section{Concepto y clasificación}

El derecho real de superficie presenta una doble naturaleza según lo dispone el CCCN.

% TODO: contenido incompleto o ambiguo en las notas originales (el apunte cita el Art. 1888 tanto para la clasificación como para la carga o gravamen real)
\begin{important}
\textbf{Art.~1888 CCCN.} El derecho real de superficie es un derecho real \textbf{sobre cosa propia} cuando existe propiedad superficiaria, y es un derecho real \textbf{sobre cosa ajena} cuando recae sobre el derecho de plantar, construir o forestar.
\end{important}

Se trata, por lo tanto, de un derecho que contempla dos vertientes bien diferenciadas, que se analizan al tratar su objeto.

\section{Caracteres generales}

Los lineamientos generales de los derechos reales se aplican al derecho de superficie de la siguiente manera:

\begin{itemize}
    \item \textbf{Derecho principal} (Art.~1889 CCCN).
    \item \textbf{Se ejerce por la posesión} (Art.~1891 CCCN).
    \item \textbf{Constituye una carga o gravamen real:} todo derecho que el titular de dominio constituya a favor de otro respecto de la cosa constituye una carga o gravamen real (Art.~1888 CCCN).
    \item \textbf{Se adquiere por título suficiente y modo suficiente} (Art.~1892 CCCN).
    \item Siempre se constituye sobre un \textbf{inmueble ajeno}.
    \item Se constituye por \textbf{escritura pública}, porque recae sobre inmuebles y el título suficiente debe contenerse en dicho instrumento.
    \item La \textbf{inscripción registral} otorga oponibilidad \textit{erga omnes} a partir de su registración.
\end{itemize}

\section{El derecho de superficie en sí mismo}

El derecho de superficie es una figura novedosa. Sus características esenciales son:

\begin{itemize}
    \item Es \textbf{temporario}.
    \item Recae sobre un \textbf{inmueble ajeno}.
    \item Confiere al titular las facultades de \textbf{uso, goce y disposición material y jurídica}.
\end{itemize}

Para precisar el alcance de la disposición conviene comparar tres situaciones. Quien es propietario tiene derecho sobre toda la cosa; quien es condómino tiene derecho sobre su parte indivisa; el superficiario, en cambio, tiene derecho a disponer y transmitir únicamente aquello que constituye el objeto de su derecho. No dispone de toda la cosa, porque la superficie siempre se constituye sobre un inmueble ajeno.

\chapter[Objeto y emplazamiento]{Objeto, emplazamiento y legitimados}

\section{Objeto del derecho de superficie}

\subsection{Una novedad del CCCN: el derecho real sobre un derecho}

Como regla general, el objeto de los derechos reales son las cosas. Solamente cuando la ley lo disponga, el derecho real puede recaer sobre un derecho. Esta es una gran novedad del CCCN respecto del Código de Vélez.

\begin{definition}{Objeto del derecho de superficie (Art.~2114 CCCN)}
El derecho de superficie puede recaer sobre el \textbf{derecho de plantar, forestar o construir} en un terreno ajeno, ya sea en el suelo, el subsuelo o el espacio aéreo.
\end{definition}

\subsection{Las dos vertientes del objeto}

\begin{definition}{Modalidades del derecho de superficie (Art.~2115 CCCN)}
El objeto del derecho puede ser:
\begin{enumerate}[label=(\alph*)]
    \item el \textbf{derecho real de construir, plantar o forestar} en un inmueble ajeno, haciendo propio aquello que se construya o plante; o
    \item la \textbf{propiedad de aquello que ya está plantado o ya está construido} en la superficie, sin adquirir la propiedad del suelo.
\end{enumerate}
\end{definition}

En todos los casos coexisten la propiedad del superficiario (ya sea el derecho real sobre el derecho a plantar, construir o forestar, o sobre lo plantado, construido o forestado) y el dominio del titular del inmueble. Este último presenta una imperfección, dada por la carga que implica la constitución de este derecho real.

\section{Emplazamiento}

El derecho puede emplazarse:

\begin{itemize}
    \item en \textbf{todo el inmueble}; o
    \item en una \textbf{parte material} del inmueble.
\end{itemize}

\section{Legitimados para constituirlo}

Pueden constituir el derecho de superficie:

\begin{itemize}
    \item únicamente el \textbf{titular de dominio};
    \item los \textbf{condóminos en su conjunto};
    \item el \textbf{propietario de un inmueble sometido a propiedad horizontal}.
\end{itemize}

Solamente ellos pueden afectar el inmueble con el gravamen que genera la constitución de un derecho real a favor de terceros, en este caso, del superficiario.

\chapter[Adquisición y facultades]{Adquisición, facultades del superficiario y situación del propietario}

\section{Adquisición del derecho}

El derecho de superficie se adquiere con \textbf{título suficiente y modo}. Además:

\begin{itemize}
    \item El acto causal de la transmisión puede ser un contrato \textbf{oneroso o gratuito}.
    \item Puede realizarse por actos \textbf{entre vivos} o \textbf{mortis causa}.
    \item \textbf{En ningún caso se adquiere por prescripción adquisitiva.}
\end{itemize}

\section{Facultades del superficiario}

Se trata de un derecho real de especial fortaleza: el superficiario, aunque no ha adquirido la propiedad del terreno, puede constituir derechos reales de garantía sobre el objeto de su derecho, es decir, sobre el derecho de construir, plantar o forestar, o sobre la propiedad superficiaria.

Las facultades reconocidas son las siguientes:

\begin{itemize}
    \item El derecho se \textbf{ejerce por la posesión} en todos los casos.
    \item Puede constituir \textbf{derechos reales de garantía} sobre el objeto de su derecho (no sobre el inmueble del propietario).
    \item Puede \textbf{hipotecarse un derecho}: la hipoteca puede recaer sobre el derecho de construir, plantar o forestar, o bien sobre la propiedad superficiaria. Esta es una particularidad remarcable.
    \item Salvo pacto en contrario en el instrumento constitutivo, puede construir un inmueble y \textbf{afectarlo a propiedad horizontal}.
    \item Puede \textbf{transmitir} lo que tiene: el derecho de superficie.
\end{itemize}

\begin{important}
Todos estos derechos están limitados en el tiempo, y esas limitaciones son de orden público.
\end{important}

\section{Situación del propietario del suelo}

Durante la vigencia del derecho de superficie, el propietario del suelo conserva la disposición material y jurídica, porque sigue siendo el dueño. Sin embargo, su dominio es imperfecto.

\begin{definition}{Dominio imperfecto del propietario}
Es la situación del titular de dominio cuyo derecho está limitado por la \textbf{prohibición de turbar el derecho del superficiario}. Debe respetarlo y, si dispone o enajena el inmueble, lo enajena con el derecho de superficie.
\end{definition}

\chapter[Duración y extinción]{Duración, destrucción del objeto y extinción}

\section{Plazos}

% TODO: contenido incompleto o ambiguo en las notas originales (el índice temático del apunte menciona «plazos máximos legales», «plazo convencional» y «orden público», pero el desarrollo no detalla los plazos máximos)
Las limitaciones temporales del derecho son de orden público. En la práctica, resulta difícil que el derecho se constituya por un plazo pequeño; en general se contempla el plazo máximo legal.

\section{Destrucción del objeto}

\begin{important}
\textbf{Art.~2122 CCCN.} El derecho real de superficie \textbf{no se extingue por la destrucción del objeto}. Si se destruyen las construcciones, el superficiario debe reconstruir en el plazo de \textbf{seis (6) años}; cuando el derecho recae sobre plantaciones o forestaciones, debe plantar o forestar en el plazo de \textbf{tres (3) años}. Si no cumple, opera la extinción del derecho.
\end{important}

En consecuencia, la extinción del objeto del derecho de superficie no produce la extinción inmediata del derecho.

\section{Modos de extinción}

El derecho de superficie se extingue por:

\begin{enumerate}
    \item \textbf{Vencimiento del plazo} convencional o legal.
    \item \textbf{Renuncia expresa} del superficiario.
    \item \textbf{Condición resolutoria}, cuando se cumple la condición a la que el derecho estaba sujeto.
    \item \textbf{Consolidación}: se da cuando las calidades de propietario y superficiario se confunden en una misma persona, de modo que el derecho de superficie se extingue.
    \item \textbf{No uso durante 10 años}, cuando el derecho fue constituido para construir.
    \item \textbf{No uso durante 5 años}, cuando el objeto fue plantar o forestar.
\end{enumerate}

\section{Relación del superficiario con terceros}

El superficiario puede constituir derechos reales y personales a favor de terceros, siempre limitados al plazo de duración de la superficie. Los efectos de la extinción sobre esos derechos dependen del modo en que aquella se produzca:

\begin{table}[H]
\centering
\renewcommand{\arraystretch}{1.3}
\begin{tabularx}{\textwidth}{
    >{\raggedright\arraybackslash}p{0.30\textwidth}
    >{\raggedright\arraybackslash}X
}
\toprule
\textbf{Modo de extinción} & \textbf{Efectos sobre terceros} \\
\midrule
Extinción anticipada (antes del plazo) &
Los derechos reales y personales constituidos a favor de terceros se mantienen y conservan sus efectos hasta la extinción del plazo original. \\
\midrule
Extinción por vencimiento del plazo &
El propietario adquiere todo lo construido, plantado o forestado que exista al vencimiento, libre de los derechos reales o personales que el superficiario haya constituido. \\
\bottomrule
\end{tabularx}
\end{table}

\chapter[Indemnización y normas supletorias]{Indemnización al superficiario y normas supletorias}

\section{Indemnización al superficiario (Art.~2126 CCCN)}

\begin{important}
Producida la extinción, el titular del suelo, que es quien adquiere lo plantado, construido o forestado, \textbf{debe indemnizar al superficiario}, excepto pacto en contrario. El monto se fija:
\begin{enumerate}[label=(\alph*)]
    \item por las partes en el acto constitutivo;
    \item por acuerdos posteriores; y
    \item si no lo hubieran fijado, por \textbf{resolución judicial}, tomando en cuenta los \textbf{valores incorporados al inmueble en los últimos dos años} y \textbf{sin descontar amortizaciones}.
\end{enumerate}
\end{important}

La aplicación práctica presenta complejidad: resulta difícil prever, con décadas de anticipación (por ejemplo, 50 o 70 años), una indemnización que contemple todo lo incorporado a lo largo del tiempo. Por esa razón se trata de un derecho que probablemente genere mucho interés en el futuro, dado que será difícil que se constituya por un plazo breve y, en general, se contempla el plazo máximo legal.

\section{Normas supletorias aplicables}

\begin{table}[H]
\centering
\renewcommand{\arraystretch}{1.3}
\begin{tabularx}{\textwidth}{
    >{\raggedright\arraybackslash}p{0.38\textwidth}
    >{\raggedright\arraybackslash}X
}
\toprule
\textbf{Tipo de superficie} & \textbf{Normas supletorias} \\
\midrule
Derecho a construir, plantar o forestar (recae sobre un derecho) &
Normas de las limitaciones de uso y goce del usufructo, en forma supletoria. \\
\midrule
Propiedad superficiaria (construcciones, forestaciones o plantaciones ya existentes) &
Normas del espíritu del instituto y, supletoriamente, las normas del dominio revocable. \\
\bottomrule
\end{tabularx}
\end{table}

\begin{note}
Dado que no existe experiencia previa con esta figura, se trata de un derecho novedoso: será necesario que transcurra el tiempo y que se abra paso a nuevas modalidades negociales que pongan en funcionamiento este derecho real.
\end{note}

% ------------------------------------------------------------
% CORNELL
% ------------------------------------------------------------
\chapter[Cuadro Cornell]{Cuadro Cornell: Derecho Real de Superficie}

\renewcommand{\arraystretch}{1.25}
\begin{longtable}{
    >{\raggedright\arraybackslash}p{0.28\textwidth}
    >{\raggedright\arraybackslash}p{0.65\textwidth}
}
\toprule
\textbf{Preguntas / palabras clave} &
\textbf{Notas / respuestas} \\
\midrule
\endfirsthead
\toprule
\textbf{Preguntas / palabras clave} &
\textbf{Notas / respuestas} \\
\midrule
\endhead
\midrule
\multicolumn{2}{r}{\small\textit{Continúa en la página siguiente}} \\
\endfoot
\endlastfoot

\textbf{¿Cómo clasifica el Art.~1888 al derecho de superficie?} &
Es un derecho real sobre cosa propia cuando existe propiedad superficiaria, y sobre cosa ajena cuando recae sobre el derecho de plantar, construir o forestar. \\[0.6em]

\textbf{¿Cuáles son sus caracteres generales?} &
Es principal (Art.~1889); se ejerce por la posesión (Art.~1891); constituye carga sobre el dominio (Art.~1888); se adquiere por título suficiente y modo (Art.~1892); se constituye sobre inmueble ajeno, por escritura pública; es oponible \textit{erga omnes} desde la inscripción registral. \\[0.6em]

\textbf{¿Qué novedad plantea el CCCN respecto del Código de Vélez?} &
El Art.~2114 permite que el objeto de un derecho real sea un derecho, y no solo una cosa. Así, el derecho de superficie puede recaer sobre el derecho de plantar, construir o forestar, como excepción a la regla general. \\[0.6em]

\textbf{¿Cuáles son las dos vertientes del objeto? (Art.~2115)} &
(a) El derecho real de construir, plantar o forestar en inmueble ajeno, haciendo propio lo que se construya o plante. (b) La propiedad de lo ya construido o plantado en la superficie, sin adquirir el suelo. \\[0.6em]

\textbf{¿Quiénes pueden constituir este derecho?} &
El titular de dominio, los condóminos en su conjunto o el propietario de un inmueble sometido a propiedad horizontal. \\[0.6em]

\textbf{¿Cómo se adquiere? ¿Puede adquirirse por prescripción?} &
Por título suficiente y modo; el acto causal puede ser oneroso o gratuito, entre vivos o \textit{mortis causa}. En ningún caso se adquiere por prescripción adquisitiva. \\[0.6em]

\textbf{¿Qué facultades tiene el superficiario?} &
Ejerce la posesión; puede constituir derechos reales de garantía sobre el objeto de su derecho; puede hipotecar el derecho de construir, plantar o forestar o la propiedad superficiaria; puede afectar lo construido a propiedad horizontal (salvo pacto en contrario) y transmitir su derecho. \\[0.6em]

\textbf{¿Qué imperfección tiene el dominio del propietario del suelo?} &
No puede turbar el ejercicio del derecho del superficiario. Si enajena el inmueble, lo transmite con la carga de la superficie. Es un dominio imperfecto durante toda la vigencia del derecho. \\[0.6em]

\textbf{¿La destrucción del objeto extingue el derecho? (Art.~2122)} &
No de forma inmediata. El superficiario debe reconstruir en 6 años (construcciones) o en 3 años (plantaciones o forestaciones). Si no lo hace, opera la extinción. \\[0.6em]

\textbf{¿Cuáles son los modos de extinción?} &
Vencimiento del plazo convencional o legal; renuncia expresa; condición resolutoria cumplida; consolidación (confusión de propietario y superficiario en una misma persona); no uso (10 años para construir, 5 años para plantar o forestar). \\[0.6em]

\textbf{¿Qué ocurre con los derechos de terceros al extinguirse la superficie?} &
Si se extingue por vencimiento del plazo, el propietario adquiere lo construido o plantado libre de cargas. Si se extingue antes del plazo, los derechos constituidos a favor de terceros subsisten hasta el vencimiento del plazo original. \\[0.6em]

\textbf{¿Cómo se fija la indemnización al superficiario? (Art.~2126)} &
Por pacto en el acto constitutivo, por acuerdos posteriores o por resolución judicial. Se consideran los valores incorporados en los últimos 2 años, sin descontar amortizaciones. Procede salvo pacto en contrario. \\[0.6em]

\textbf{¿Qué normas supletorias aplican?} &
Al derecho sobre un derecho (plantar, construir o forestar): normas de las limitaciones de uso y goce del usufructo. A la propiedad superficiaria (sobre cosas ya existentes): normas del dominio revocable. \\

\midrule
\multicolumn{2}{p{0.93\textwidth}}{
\textbf{Resumen:} El derecho real de superficie, regulado en los Arts.~2114 a 2126 del CCCN, se caracteriza por la disociación entre la propiedad del suelo y la propiedad de lo construido, plantado o forestado sobre él: el superficiario puede ser dueño de lo edificado sin ser dueño del terreno. Admite dos vertientes según su objeto: (a) el derecho real de construir, plantar o forestar en terreno ajeno, caso en que el objeto es, excepcionalmente, un derecho y no una cosa, y (b) la propiedad superficiaria sobre lo ya construido o plantado. Se adquiere por título suficiente y modo (nunca por prescripción), requiere escritura pública e inscripción registral para su oponibilidad \textit{erga omnes}, y solo pueden constituirlo el titular de dominio, los condóminos en conjunto o el propietario en propiedad horizontal. El superficiario ejerce la posesión, puede constituir derechos reales de garantía, hipotecar el derecho mismo y afectar lo construido a propiedad horizontal. El dominio del propietario queda gravado con la prohibición de turbar el ejercicio del superficiario. La extinción opera por vencimiento del plazo, renuncia, condición resolutoria, consolidación o no uso (10 años para construcciones; 5 para plantaciones o forestaciones); la destrucción del objeto no la produce automáticamente, pues el superficiario dispone de 6 o 3 años para reconstruir o replantar. Extinguido el derecho, el propietario adquiere lo construido o plantado y debe indemnizar al superficiario, salvo pacto en contrario, sobre la base de los valores incorporados en los últimos dos años y sin descontar amortizaciones. Supletoriamente rigen las normas del usufructo (limitaciones de uso y goce) y las del dominio revocable (propiedad superficiaria).
} \\
\bottomrule
\end{longtable}

\end{document}
